\chapter{年表}

\textbf{1940} 阴历七月初九出生于石桥街西夹后布袋街外祖母家的巷子里。辗转生活于薛庄、沙山、石人沟、魏庄。

\textbf{1941} 阴历五月十七妻子王若平出生于蒲山镇。

\textbf{1948} 患天花

\textbf{1949} 新中国成立

\textbf{1952} 土改，在魏庄分得三间大瓦房和几亩地

\textbf{1956} 读初中

\textbf{1959} 秋季，就读于郑州农业机械化专科学校（一年后更名为河南农学院农机分院），现为河南工学院。

\textbf{1961} 全国大中院校放长假，停课回家，在生产队种地。

\textbf{1962} 农历七月十一与王若平结婚，住在魏庄

\textbf{1963} 阴历五月初一，大儿子陈华军出生于魏庄，送至沽沱抚养。八月下旬与妻子重回郑州上学，九月一日报道。

\textbf{1964} 八月中旬毕业分配在茶庵工作，妻子分配至安皋，两地分居。

\textbf{1966} 四月，调动至地方国营南阳县拖拉机总站机务科（农机局前 身）

\textbf{1967} 阴历五月二十九，二儿子陈永出生于南阳地区中心医院。阴 历六月十三，儿媳杜思桂出生

\textbf{1968} 秋天，妻子担任拖拉机修配厂车工。

\textbf{1969} 阴历四月初三，三儿子陈斌于南阳市医院出生。

\textbf{1971} 调动至农业局工作，在李华庄药厂担任政工组长。阴历十一月十七，小儿子陈涛于南阳地区中心医院出生。

\textbf{1972} 在煤厂街盖房子，由妻子主要负责。

\textbf{1974} 妻子调动至南阳县计委仓库保管员，一边工作一边照顾四个儿子。

\textbf{1976} 调动至县委办公室，工作几个月时间。

\textbf{1976} 阴历十月十九，儿媳宋爱平出生。

\textbf{1977} 调至组织部工作。

\textbf{1978} 阴历十月初一，儿媳刘延娇出生。

\textbf{1980} 阴历三月二十九，儿媳张娟出生。

\textbf{1980} 阴历七月十二父亲陈锡海逝世，葬于皇路店魏庄

\textbf{1983} 调动至县人事局右派改正办公室工作，主要负责右派改正。

\textbf{1984} 调动至南阳县农机局监理站工作。

\textbf{1985} 腊月初一母亲吕金兰逝世，葬于皇路店魏庄

\textbf{1989} 阴历八八年腊月十九陈华军与范惠杰结婚、腊月二十陈永与 杜思桂结婚。

\textbf{1989} 阴历九月初八孙女陈旭鲲出生于南阳地区中心医院。

\textbf{1996} 妻子退休，55岁。单位为县物资局金属公司。

\textbf{1995} 阴历八月二十五孙子陈旭鹏出生于南阳地区中心医院。

\textbf{1999} 阴历二月二十六陈斌与宋爱平结婚。

\textbf{2000} 因身体原因退休，60岁。

\textbf{2001} 阴历正月初七孙女陈泓妤出生于南阳市县医院。

\textbf{2002} 阴历十一月初九陈涛与刘延娇结婚。

\textbf{2003} 阴历六月初二孙女陈奕霏出生于南阳市县医院。

\textbf{2004} 阴历三月初九陈华军与张娟结婚。 阴历十一月初四孙女陈佳欣出生于南阳市县医院。

\textbf{2012} 阴历五月十二孙女陈映竹出生于南阳地区中心医院。

\textbf{2016} 阴历一五年腊月十三孙子陈旭宏出生于南阳市县医院。

\textbf{2017} 阴历三月十九，重建老家魏庄的住宅。

\textbf{2017} 阴历九月初三孙女陈旭鲲与刘一佳结婚。

\textbf{2017} 年末，老家新房建成。

\textbf{2018} 阴历六月十八重孙女刘沁孜出生。

\textbf{2025} 5月25日，孙子陈旭鹏与孟舒晨结婚。

\textbf{2026} 5月24日15时28分，在南阳市中心医院安然离世，享年八十七岁。


\clearpage
\begin{landscape}
\thispagestyle{plain}
\null\vspace*{-0.4cm}
\begin{center}
\begin{tikzpicture}[
  remember picture,
  every node/.style={
    draw=seal, rounded corners=2.5pt, line width=0.45pt,
    align=center, font=\trueka\small,
    inner sep=3.5pt, minimum width=1.55cm
  },
  anc/.style={
    fill=seal!6, font=\trueka\small\bfseries,
    minimum width=1.85cm, inner sep=4.5pt
  },
  spouse/.style={
    fill=seal!10, font=\trueka\small\bfseries,
    minimum width=1.85cm, inner sep=4.5pt
  },
  small/.style={font=\trueka\footnotesize, minimum width=1.35cm, inner sep=2.8pt},
  ln/.style={draw=seal!60, line width=0.45pt},
  marriage/.style={draw=seal!75, line width=0.7pt, double, double distance=1pt},
  x=1cm, y=1cm,
]

% ── Columns ────────────────────────────────────────────
\def\cA{0}      % 祖辈
\def\cB{3.2}    % 炳林一代
\def\cC{6.4}    % 子女
\def\cD{9.5}    % 孙辈
\def\cE{12.6}   % 重孙辈

% ── 标题（嵌入 tikz 内部，确保与图同页） ─────────────
\node[draw=none, font=\trueka\fontsize{18pt}{22pt}\selectfont\bfseries, text=inkblack]
  (title) at (\cC, 7.2) {家\hspace{0.8em}族\hspace{0.8em}世\hspace{0.8em}系};
\node[draw=none, below=2pt of title]
  {{\color{seal}\rule{1.2em}{0.4pt}}\hspace{0.4em}%
    \raisebox{-1pt}{\color{seal}\small$\Diamond$}\hspace{0.4em}%
    {\color{seal}\rule{1.2em}{0.4pt}}};

% ── 第一代：双方父母 ──────────────────────────────────
\node[anc] (chenP) at (\cA, 1.6)  {陈\,锡\,海\\吕\,金\,兰};
\node[anc] (wangP) at (\cA, -2.3) {王\,铭\,鼎\\庆\hspace{1em}良};

% ── 第二代：炳林兄妹 + 若平 ──────────────────────────
\node (cby)  at (\cB,  4.2) {陈炳郁};
\node (cbq)  at (\cB,  3.3) {陈炳琴};
\node (cbi)  at (\cB,  2.4) {陈炳义};
\node (cbll) at (\cB,  1.5) {陈炳立};
\node (cbyu) at (\cB,  0.6) {陈炳玉};
\node[spouse] (cbl) at (\cB, -0.65) {陈\,炳\,林};
\node[spouse] (wrp) at (\cB, -2.30) {王\,若\,平};

% 婚姻线（双线表示婚配）
\draw[marriage] (cbl.south) -- (wrp.north);
\coordinate (union) at ($(cbl.south)!0.5!(wrp.north)$);

% 陈家父母 → 五位兄妹 + 炳林
\foreach \n in {cby,cbq,cbi,cbll,cbyu,cbl} {
  \draw[ln] (chenP.east) to[out=0,in=180] (\n.west);
}
% 王家父母 → 若平
\draw[ln] (wangP.east) to[out=0,in=180] (wrp.west);

% ── 第三代：炳林若平之子女 ────────────────────────────
\node (chj) at (\cC,  3.6) {陈\,华\,军};
\node (cy)  at (\cC,  1.2) {陈\hspace{1em}永\\杜思桂};
\node (cb)  at (\cC, -1.2) {陈\hspace{1em}斌\\宋爱平};
\node (ct)  at (\cC, -3.6) {陈\hspace{1em}涛\\刘延娇};

% 由婚姻线汇流分枝
\coordinate (junC) at (\cC-1.4, -1.2);
\draw[ln] (union) -- (junC);
\foreach \n in {chj,cy,cb,ct} {
  \draw[ln] (junC) to[out=0,in=180] (\n.west);
}

% 华军两段婚姻：前妻 / 张娟
\node[small] (qq) at (\cC+1.7,  4.3) {（前妻）};
\node[small] (zj) at (\cC+1.7,  2.9) {张\hspace{0.6em}娟};
\draw[ln] (chj.east) to[out=0,in=180] (qq.west);
\draw[ln] (chj.east) to[out=0,in=180] (zj.west);

% ── 第四代：孙辈 ─────────────────────────────────────
\node (xk) at (\cD,  4.7) {陈旭鲲\\刘一佳};
\node (jx) at (\cD,  2.9) {陈佳欣};
\node (xp) at (\cD,  1.2) {陈旭鹏\\孟舒晨};
\node (hy) at (\cD, -0.5) {陈泓妤};
\node (yz) at (\cD, -1.9) {陈映竹};
\node (yf) at (\cD, -2.9) {陈奕霏};
\node (xh) at (\cD, -4.3) {陈旭宏};

\draw[ln] (qq.east) to[out=0,in=180] (xk.west);
\draw[ln] (zj.east) to[out=0,in=180] (jx.west);
\draw[ln] (cy.east) to[out=0,in=180] (xp.west);
\draw[ln] (cb.east) to[out=0,in=180] (hy.west);
\draw[ln] (cb.east) to[out=0,in=180] (yz.west);
\draw[ln] (ct.east) to[out=0,in=180] (yf.west);
\draw[ln] (ct.east) to[out=0,in=180] (xh.west);

% ── 第五代：重孙辈 ──────────────────────────────────
\node (qz) at (\cE, 4.7) {刘沁孜};
\draw[ln] (xk.east) to[out=0,in=180] (qz.west);

% ── 代际题注（横排列于上方） ────────────────────────
\node[draw=none, font=\trueka\footnotesize\itshape, text=inkgray]
  at (\cA, 5.7) {祖\,辈};
\node[draw=none, font=\trueka\footnotesize\itshape, text=inkgray]
  at (\cB, 5.7) {兄\,妹};
\node[draw=none, font=\trueka\footnotesize\itshape, text=inkgray]
  at (\cC, 5.7) {子\,女};
\node[draw=none, font=\trueka\footnotesize\itshape, text=inkgray]
  at (\cD, 5.7) {孙\,辈};
\node[draw=none, font=\trueka\footnotesize\itshape, text=inkgray]
  at (\cE, 5.7) {重\,孙};

\end{tikzpicture}
\end{center}
\end{landscape}

